% Generated draft: requires mathematical and source review.
\documentclass[11pt]{article}
\usepackage[margin=25mm]{geometry}
\usepackage{fontspec}
\setmainfont{Latin Modern Roman}
\usepackage{amsmath,amssymb,mathtools,unicode-math}
\setmathfont{Latin Modern Math}
\usepackage{enumitem,needspace,xurl,hyperref}
\setlength{\parindent}{0pt}
\setlength{\parskip}{5pt}
\setlength{\emergencystretch}{3em}
\newcommand{\N}{\mathbb{N}}
\newcommand{\Z}{\mathbb{Z}}
\newcommand{\Q}{\mathbb{Q}}
\newcommand{\R}{\mathbb{R}}
\newcommand{\C}{\mathbb{C}}
\newcommand{\F}{\mathbb{F}}
\newcommand{\A}{\mathbb{A}}
\newcommand{\PP}{\mathbb{P}}
\newcommand{\E}{\mathbb{E}}
\newcommand{\eps}{\varepsilon}
\newcommand{\dd}{\,\mathrm{d}}
\begin{document}
\section*{UnsolvedMath ID: 20001935}
\textbf{Source problem code:} catalogue entry
\subsection{Definitions and mathematical statement}
Let \(M\) be a closed smooth manifold of positive dimension. For a Riemannian metric \(g\), let \(S_gM\) be its unit tangent bundle, \(\phi_g^t\) its geodesic flow, and
\[
\operatorname{Per}(g)=\{v\in S_gM:\phi_g^T(v)=v\text{ for some }T>0\}.
\]
Question 11.2.2 asks whether \(\overline{\operatorname{Per}(g)}=S_gM\) for a generic metric. This is density of directions in phase space, not just density of closed-geodesic images in \(M\). Allowing positive rescalings of periodic unit vectors gives the corresponding nonzero-tangent-bundle statement.

The source does not specify the genericity topology. For a precise residual version use the Baire space \(\operatorname{Met}^r(M)\) of \(C^r\) metrics, with its \(C^r\) topology, for an integer \(r\geq2\); alternatively use smooth metrics with the \(C^\infty\) topology. The assertion is that the metrics with dense periodic directions contain a countable intersection of open dense subsets. Smooth metrics equipped only with a weaker finite \(C^r\) topology must not silently be treated as a complete Baire model.

The supplied aid reports a conditional phase-box reduction. Fix an auxiliary smooth metric \(h\), a countable basis \((U_j)\) of nonempty open sets in \(S_hM\), and identifications \(I_g(v)=v/\sqrt{g(v,v)}\). Suppose every metric neighborhood contains a metric with a nondegenerate periodic geodesic meeting \(I_g(U_j)\), for every \(j\). Nondegenerate means no eigenvalue one in the transverse Poincar\'e return map. Persistence makes the corresponding box-hitting sets open; the stated creation property makes them dense. Their intersection gives residual density of periodic directions. This is a reduction to localized metric perturbations, not a proof that such perturbations always exist.

The aid also reports density when Liouville measure has full support and is nonuniformly hyperbolic, meaning all Lyapunov exponents transverse to the flow direction are nonzero almost everywhere. Recurrence and the Pesin--Katok closing lemma are the stated ingredients. Its literature audit distinguishes Rifford's \(C^1\) metric closing result, Bessa--Torres's residual theorem in a \(C^0\) closure of flows, Riemannian surface base-density/equidistribution results, and Chen's phase-density theorem for reversible Finsler surfaces. It says these do not by themselves settle the original higher-regularity Riemannian phase-density problem. Full theorem references are not supplied.

The large appended section 12 concerns a separate reconstruction problem. It defines a smooth path structure with one unparameterized curve through each tangent line. For a torsion-free connection, affine parameters require
\[
\ddot\gamma^i+\Gamma^i_{jk}\dot\gamma^j\dot\gamma^k=0,
\]
while arbitrary parameters permit a curve-dependent scalar multiple of \(\dot\gamma^i\) on the right. There are \(n^2(n+1)/2\) symmetric Christoffel unknowns at a point; affine data determine them uniquely when the velocity dyads span \(\operatorname{Sym}^2(T_pM)\), but all other curves still require consistency. Unparameterized data determine only a projective class,
\[
\overline\Gamma^i_{jk}=\Gamma^i_{jk}+\delta^i_j\varphi_k+\delta^i_k\varphi_j.
\]
A metric with that affine connection requires a nondegenerate symmetric form preserved by full holonomy. The source discusses linear curvature-symmetry tests and all covariant derivatives, analytic versus smooth integrability limitations, local conformal/parallel-volume reconstruction, and local classifications of metrics sharing one Levi-Civita connection, with a Riemannian global de Rham counterpart. Global volume reconstruction needs more than a local closedness test. The extract ends before its announced next unsolved question. None of these appended reconstruction claims supplies a solution to periodic phase density.
\subsection{Short English statement}
Are periodic geodesic directions dense for a generic Riemannian metric? The supplied aid reduces a precise residual version to nondegenerate orbit creation in every phase-space box and gives the nonuniformly hyperbolic case. Lower-regularity, base-density and larger Finsler-class results are kept distinct.
\subsection{Sources}
\begin{itemize}
\item Keith Burns and Vladimir Matveev, Question 11.2.2, AIM geodesics problem list: \url{https://aimath.org/pastworkshops/geodesicsproblems.pdf}.
\item The aid names Rifford, Bessa--Torres, Chen and the Pesin--Katok closing lemma without full bibliographic entries.
\item Appended section 12 cites [Thomas1925], [Matveev2012a, section 2.1], [Levi-Civita1896], [Weyl1921], Bryant and Thurston [Bry-Thu2011], [Schmidt1973], [Mat-Trau2014], Cartan and Eisenhart [Eisenhart1923], [Boubel2012, Boubel2014], and [DeRham1952]. Full entries are not supplied.
\item Local registry \texttt{data/unsolvedmath/problems.json}, permanent ID 20001935, code \texttt{catalogue entry}; canonical \texttt{aim-geometry-notes.json}, \texttt{notes[272]}, source item 11.2; updated 14 August 2026.
\item Attempt 1: reduction, problem status at run open, rigor label proved, novelty confidence low; current registry status partially solved. The machine-generated aid is not peer reviewed and requires independent verification before citation.
\end{itemize}
\end{document}
